\subsection{Short-Production Grammar Normalization and Consuming Universality}
\label{annex:p09-short-grammar-normalization}

\begin{annexitem}{annex:p09.short-grammar-normalization}{thm:p09-short-grammar-normalization}{Finite compiler, unique control, and exact language preservation}
\begin{proof}
\paragraph{Effective construction and production bounds.}
The finite input provides all letters of \(N\cup T\), finitely many
productions, and the finite lengths of their sides. Assign distinct labels
to those productions and choose a finite supply of new names disjoint from
the input alphabets. The listed loops enumerate the coded letters,
initialization and ready controls, and all matching/emission states. They
then enumerate all ten finite production groups. This terminates and
constructs a finite explicit grammar, not a semantic existence selection.

Initialization and emission use one-to-two productions; commitment, mode
change and ready restoration use one-to-one productions; cursor motion and
terminal emission use two-to-two productions; matching uses two-to-one
productions; cleanup start uses two-to-one and cleanup finish two-to-zero.
Thus every left side has length one or two and every right side length
zero, one or two. All left sides contain a fresh nonterminal.
Every one-to-one production changes to a distinct control by freshness.
For cursor motion, the cursor and coded letter are distinct, so the two
words differ in order. Every other production either changes length or
introduces a distinct control or an original terminal. Hence none is an
identity. The terminal alphabet remains exactly \(T\).

\paragraph{Initialization cannot interleave with simulation.}
From \(S^\dagger\), the only derivation prefix is
\[
 S^\dagger\Longrightarrow LI
 \Longrightarrow LA_SJ
 \Longrightarrow LA_SQR.
\]
Before the last step no \(Q\), matching state, emission state or \(D\)
exists, so none of their productions can apply. Coded symbols have no
standalone productions. The only available next left side is the indicated
initialization control. After the final step, \(I,J,S^\dagger\) are absent
and no production can recreate them. The first ready word decodes to \(S\).

\paragraph{Unique-control structure.}
Once initialization finishes and before a completely terminal word is
obtained, every reachable word contains exactly one control from
\[
 \{Q,D\}\cup\{Q_{r,i}\}\cup\{E_{r,j}\}.
\]
Every production in modes other than initialization has precisely one of
these controls on its left, and replaces it by precisely one control,
except final cleanup, which removes \(D\) and \(R\).
No coded data symbol or original terminal can initiate a production alone.
Induction on derivation length proves this uniqueness and also ensures that
productions from two different modes cannot operate at separate positions.

More explicitly, all reachable words after initialization lie in the
following classes, with \(p,q\in(N\cup T)^*\) and \(z\in T^*\):
\begin{enumerate}
\item ready: \(L\widehat pQ\widehat qR\);
\item matching: \(L\widehat pQ_{r,i}\widehat qR\);
\item emitting: \(L\widehat pE_{r,j}\widehat qR\);
\item cleanup: \(zD\widehat qR\);
\item terminal: \(z\).
\end{enumerate}
The matching and emitting classes are constrained by their entry histories,
as proved next; this list is not an assumption that an arbitrary word of
one of those shapes is a correct simulation. The boundary \(L\) remains
unique and leftmost in the first three classes, and \(R\) unique and
rightmost in the first four. Cursors cannot cross either boundary: motion
rules match only a coded data symbol, never \(L\) or \(R\).

\paragraph{Ready cursor motion.}
In a ready word, a move right replaces \(QA_a\) by \(A_aQ\), transferring
one decoded letter from \(q\) to \(p\). A move left does the converse.
In either case the concatenation \(pq\) is unchanged. Repeated moves
therefore place the unique ready cursor at every gap of the finite decoded
word, including both ends. None changes a coded letter or a boundary.

The other productions available in ready mode are a commitment
\(Q\to Q_{r,0}\), or \(LQ\to D\) when the left boundary is immediately
adjacent. No emission or match-state production can operate while \(Q\)
is the control. Thus the stated alternatives exhaust all ready steps.

\paragraph{Matching and invalid-match behavior.}
Suppose commitment occurs in \(L\widehat pQ\widehat qR\), selecting
\(r:u_r=a_{r,1}\cdots a_{r,m_r}\to v_r\).
At state \(Q_{r,0}\), the only production with that control on its left
is \(Q_{r,0}A_{a_{r,1}}\to Q_{r,1}\).
For each \(0\leq i<m_r\), the only left side with \(Q_{r,i}\) is the
matching pair for \(a_{r,i+1}\). It cannot switch production label,
move the cursor, restore \(Q\), or start cleanup.

Induction on \(i\) shows that reaching \(Q_{r,i}\) is possible exactly
when \(q\) has prefix \(a_{r,1}\cdots a_{r,i}\); the reached word is
\[
 L\widehat pQ_{r,i}\widehat tR,
 \qquad q=a_{r,1}\cdots a_{r,i}t.
\]
At the next mismatch, \(\widehat t\) begins with a different coded letter
or is empty so that \(R\) follows. In either case the sole matching left
side does not occur. Every other production requires a different control,
an absent initialization symbol, or \(D\); no coded symbol has a standalone
production. The whole word is therefore stuck, not just locally blocked.
It contains a nonterminal control and cannot be terminal.

If all \(m_r\) letters match, the only available continuation is
\(Q_{r,m_r}\to E_{r,0}\). It is taken at
\(L\widehat pQ_{r,m_r}\widehat tR\), with the actual precommitment
decoding \(pu_rt\). No shortened left side or speculative unmatched
prefix can reach emission mode.

\paragraph{Emission and the completed macrostep.}
For \(v_r=b_{r,1}\cdots b_{r,n_r}\), emission state \(E_{r,j}\),
\(j<n_r\), has exactly the one production
\(E_{r,j}\to A_{b_{r,j+1}}E_{r,j+1}\).
It appends that next coded right-side letter immediately to the left
of the control. Induction on \(j\) gives precisely
\[
 L\widehat{pb_{r,1}\cdots b_{r,j}}E_{r,j}\widehat tR.
\]
No match or movement production can interleave: its distinct control is
absent. At \(j=n_r\), only \(E_{r,n_r}\to Q\) is available, giving
\(L\widehat{pv_r}Q\widehat tR\).
Its decoding is \(pv_rt\), obtained from \(pu_rt\) by exactly the
chosen original grammar production in its original left/right contexts.

If \(v_r=\lambda\), then \(n_r=0\), emission mode has no insertion
step, and ready restoration gives \(L\widehat pQ\widehat tR\).
If \(u_r\) has length one, there is still one matching step.
Thus neither erasure nor the shortest nonempty left side is hidden in a
special unlisted rule.

It follows inductively that every reachable ready word decodes to a
sentential form of \(G\): the initial decoding is \(S\), movement
preserves it, and every successful excursion through match/emission modes
is one actual grammar step. Failed commitments cannot return to ready mode.

\paragraph{Every input grammar derivation is simulated.}
Given a derivation step \(pu_rt\Longrightarrow_P pv_rt\), start with
any ready word decoding to \(pu_rt\). Move \(Q\) to the gap just after
the prefix \(p\). Commit to \(r\), consume its actual matching letters,
emit its actual right side, and restore \(Q\).
The resulting ready word decodes to \(pv_rt\).
This construction is valid even when the selected occurrence is at either
word boundary or when \(t\) or \(v_r\) is empty.
Induction on an arbitrary finite grammar derivation therefore gives a ready
word for every sentential form of \(G\), and in particular for every
\(w\in L(G)\).

\paragraph{Cleanup, exact terminal language, and no symbol leakage.}
Cleanup can start only at an occurrence of \(LQ\).
In a ready word \(L\widehat pQ\widehat qR\) this requires
\(\widehat p=\lambda\), hence \(p=\lambda\) because the coding is
letter-to-letter and injective. Thus cleanup enters \(D\widehat wR\),
where the entire \(w\) is a previously established sentential form of \(G\).

If the first coded letter is \(A_t\) for \(t\in T\), the unique
applicable production is \(DA_t\to tD\). It emits exactly that original
terminal to the left of the unique cleanup control. Induction gives
\(zD\widehat qR\) with \(w=zq\), \(z\in T^*\).
If \(q\) begins with a nonterminal \(a\in N\), there is no \(DA_a\)
production. No ready/match/emission control exists and no other production
can apply, so cleanup is permanently stuck. It cannot skip that letter,
erase it, or resume grammar simulation after emitting a terminal prefix.

If \(q=\lambda\), the only remaining production is \(DR\to\lambda\).
It removes exactly the cleanup control and right boundary, leaving \(z=w\).
Consequently cleanup succeeds if and only if the whole decoded sentential
form lies in \(T^*\), and its output is that very word.
No initialization, ready, matching, or emitting word is terminal, because
each contains a fresh control and boundaries. No failed matching or failed
cleanup word is terminal, because its fresh control persists.
Every terminal word of \(G^\dagger\) therefore belongs to \(L(G)\).

Conversely, for \(w\in L(G)\), the preceding simulation produces a
ready word decoding to \(w\). Move its cursor to the left boundary,
start cleanup, emit each actual terminal of \(w\), and finish cleanup.
This gives \(w\in L(G^\dagger)\).
For \(w=\lambda\), the ready word is exactly \(LQR\); cleanup takes
\(LQR\Longrightarrow DR\Longrightarrow\lambda\).
No empty-word convention or erasing observation map is used. Both literal
inclusions, the production bounds and effective construction are proved.
\end{proof}
\end{annexitem}

\begin{annexitem}{annex:p09.consuming-universality-short-compiler}{thm:p09-consuming-universality-short-compiler}{Composition with the repaired bounded consuming simulation}
\begin{proof}
Apply Theorem~\ref{thm:p09-short-grammar-normalization} effectively to the
given finite type-0 grammar \(G\). It supplies the finite grammar
\(G^\dagger=(N^\dagger,T,S^\dagger,P^\dagger)\) with the same literal
terminal language. All its productions have nonempty left side of length
at most two and right side of length at most two; none is an identity.
Thus it satisfies precisely the grammar hypotheses in
Definition~\ref{def:p09-recombination-consuming-simulation}.

Enumerate \(P^\dagger\) with distinct production labels, use its finite
working alphabet \(N^\dagger\cup T\), and choose the fresh markers
specified in that definition disjoint from this larger alphabet.
Its bounded finite ranges enumerate a finite rule set, including the
exact right-boundary group-5 case. Construct the finite-support multiset
with one active start word \(X_1^2YS^\dagger X_2^2\), and infinite
copies of the finitely many specified reservoir words.
Every choice is a finite computation on \(G^\dagger\).

Theorem~\ref{thm:p09-recombination-consuming-simulation}, including its
actual phase invariant and terminal cleanup, now gives
\[
 \sigma^*(M)\cap T^*=L(G^\dagger)=L(G).
\]
The terminal filter is unchanged: encoded symbols and compiler controls
belong to \(N^\dagger\), not to \(T\). They cannot count as accepted
outputs simply because they occur in the full splicing language.
Both constituent equalities retain \(\lambda\), so their composition does
too. The bounded simulation's handle inspection gives radius at most five,
its support is finite, and its multiplicity-one active computation is
distinguished from the infinite reservoir supplies.

Finally, under the type-0-grammar characterization of recursively enumerable
languages, each such language has some finite grammar \(G\).
The proved effective construction applies to that grammar and gives the
asserted representation. This last use of the grammar characterization
does not claim to decide whether an arbitrary infinite language is recursively
enumerable or to recover a grammar without its finite presentation.
\end{proof}
\end{annexitem}
